\begin{abstract}
\normalfont
在消费级设备上部署多个 AI 服务，同时受到紧张的 GPU 显存预算和不同工作负载行为差异的双重约束。在我们的目标配置中，一块 8GB 显存的 GPU 无法同时从容承载一个常驻的 LLM 运行时和一个突发式图像生成运行时，而现有方案往往要么面向数据中心，要么仅针对单一模型做孤立优化。

本文提出 EdgeSched，这是一种面向本地 AI 运行时的调度器，用于在具有不同服务要求的工作负载之间复用稀缺的加速器内存。其核心机制是一种不对称迁移策略，将快速且内存安全的下行迁移，与低中断的上行恢复分离处理。系统还进一步集成了工作负载感知分发、反压控制、取消语义和崩溃恢复，以处理这一场景中的主要生命周期事件。

在 RTX 5070 Laptop（8GB VRAM）上的实验表明，在所评估的混合工作负载场景中，EdgeSched 能在突发任务执行期间以降级形式保持常驻服务可用，并能以可预测方式恢复高性能服务；相对于基于原始信号量的实现，它只引入了适度的调度开销。
\end{abstract}

\noindent\textbf{关键词：}本地 AI 运行时，资源调度，大语言模型，图像生成，GPU 内存管理